\documentclass[10pt,a4paper]{amsart}
\usepackage[margin=19mm]{geometry}
\usepackage{amsmath,amssymb,mathtools}
\usepackage[scr=rsfs,cal=euler]{mathalfa}
\usepackage{xfrac}
\usepackage[unicode,hidelinks,bookmarks=false]{hyperref}
\newcommand{\ie}{i.e.\ }
\newcommand{\cf}{cf.\ }
\newcommand{\resp}{respectively\ }
\newcommand{\Subsection}{\subsection}
\providecommand{\nobreakdash}{\nobreak\hbox{-}}
\setlength{\emergencystretch}{2em}
\multlinegap0pt
\usepackage{bm}
\usepackage{bbm}
\usepackage{dsfont}
\usepackage{xspace}
\usepackage{cancel}
\usepackage{mathtools}
\usepackage{amsthm}
\makeatletter
\@ifundefined{theorem}{\newtheorem{theorem}{Theorem}}{}
\makeatother
\title[Perspective Central Triangles Formed from a Triangle and a T]{Perspective Central Triangles Formed from a Triangle and a Transversal}
\author{Stanley Rabinowitz, Ercole Suppa}
\date{Source: arXiv:2607.25730v1 (CC BY 4.0)}
\begin{document}
\maketitle

\noindent\emph{Source and licence.} Perspective Central Triangles Formed from a Triangle and a Transversal. Version arXiv:2607.25730v1 (CC BY 4.0), distributed under the Creative Commons Attribution 4.0 International licence, as recorded in the arXiv licence metadata for this exact version and shown on the abstract page https://arxiv.org/abs/2607.25730v1. Full rights notice: licenses/arxiv-2607.25730-CC-BY-4.0.txt.\par\smallskip

\noindent\textit{Editorial scope.} Counted unit: the Theorem whose source label is \texttt{thm:x131euler} in the article (source0.tex lines 2447--2449, source label \texttt{thm:x131euler}), delivered together with the article's own complete original proof of it (source0.tex lines 2451--2583, one \texttt{proof} environment of 133 lines). No mathematics is written, repaired or re-derived: statement and proof are the article's own text line for line. Required context retained uncounted: none. Every definition, display and label of those lines is retained unchanged. Omitted from this package and not redistributed: 1--2446, 2450--2450, 2584--2691. Nothing inside a retained excerpt is omitted or altered: every retained body is delivered exactly as the article's lines are written, so no omission receipt of any kind accompanies this package. Citations: the retained excerpts carry no citation command at all, so no bibliography entry of the article is transcribed and no reference is redistributed. Reference adaptation: none was needed and none was made; every reference inside the delivered unit resolves inside it, so no unresolved reference or citation is produced. Printed slips kept literal: no printed word, symbol, hypothesis or number of the counted statement or of its proof was altered. Layout and class: the article's own class is \texttt{amsart} with packages including inputenc, babel, amsmath, amsfonts, amssymb, color, amsthm, graphics, epsfig, url, hyperref, geometry; the wrapper is the lane's pinned \texttt{amsart} editorial wrapper at 10pt on A4, which restates the theorem-like environments the retained text needs and adds the article's own macro definitions that the retained text needs (none), with extra wrapper packages (bm, bbm, dsfont, xspace, cancel, mathtools). The delivered numbering is the unit's own. Assets: no retained span contains a figure environment, a graphics-inclusion command, a PDF-page inclusion, a table or a TikZ picture; no image file is copied into this package, the package declares no asset dependency and no image file is redistributed with it. Licence: version arXiv:2607.25730v1 of this article is distributed under the Creative Commons Attribution 4.0 International licence, as recorded in the arXiv licence metadata for this exact version and shown on the abstract page https://arxiv.org/abs/2607.25730v1. A full rights notice is delivered at licenses/arxiv-2607.25730-CC-BY-4.0.txt. Characters: every retained excerpt is pure ASCII, so the delivered engine, XeLaTeX with its default Latin Modern text font, reports no missing character.
\begin{theorem}\label{thm:x131euler}
The center $X_{131}$ is concurrent for the Euler line.
\end{theorem}

\begin{proof}
To avoid conflict with the notation $S=s_As_Bs_C$ and $T=t_At_Bt_C$ of
Theorem~3.1, write
\[
\sigma=\sin 2A+\sin 2B+\sin 2C,
\qquad
\tau=\tan 2A+\tan 2B+\tan 2C.
\]
Since ETC gives the trilinear first coordinate
$(\sec A)\bigl(2\tau-\sigma(\sec 2B+\sec 2C)\bigr)\bigl(\tau-\sigma\sec 2A\bigr)$,
the corresponding barycentric center function is
\begin{equation}\label{eq:f131etc}
f_{131}
=\tan A\,\bigl(2\tau-\sigma(\sec 2B+\sec 2C)\bigr)\bigl(\tau-\sigma\sec 2A\bigr).
\end{equation}

\emph{Step 1: an algebraic form of the center function.}
Consider a triangle with side lengths $x$, $y$, $z$, angles $\alpha$,
$\beta$, $\gamma$ opposite them, and area $K$, so that
\[
16K^2=2x^2y^2+2y^2z^2+2z^2x^2-x^4-y^4-z^4 .
\]
In analogy with the Conway notation, put
\[
S_x=y^2+z^2-x^2,\qquad S_y=z^2+x^2-y^2,\qquad S_z=x^2+y^2-z^2,
\]
and abbreviate
\[
m=x^2y^2z^2,\qquad h=8K^2,\qquad
f_1=y^2z^2-h,\quad f_2=z^2x^2-h,\quad f_3=x^2y^2-h .
\]
From $\sin\alpha=2K/(yz)$ and $\cos\alpha=S_x/(2yz)$ we obtain
\[
\tan\alpha=\frac{4K}{S_x},\qquad
\sin 2\alpha=\frac{2K\,S_x}{y^2z^2},\qquad
\cos 2\alpha=1-\frac{8K^2}{y^2z^2}=\frac{f_1}{y^2z^2},
\]
and cyclically.  Hence
\[
\sigma=\frac{2K}{m}\,N_\sigma,
\qquad
N_\sigma=x^2S_x+y^2S_y+z^2S_z,
\]
\[
\tau=\frac{2K}{f_1f_2f_3}\,N_\tau,
\qquad
N_\tau=S_xf_2f_3+S_yf_1f_3+S_zf_1f_2,
\]
\[
\sec 2\alpha=\frac{y^2z^2}{f_1},
\qquad
\sec 2\beta=\frac{z^2x^2}{f_2},
\qquad
\sec 2\gamma=\frac{x^2y^2}{f_3}.
\]
Substituting these expressions into \eqref{eq:f131etc} and combining
each factor over a common denominator gives
\begin{equation}\label{eq:f131alg}
f_{131}
=\frac{16K^3\,N_1N_2}{S_x\,\bigl(mf_1f_2f_3\bigr)^{2}},
\end{equation}
where
\begin{align}
N_1&=2N_\tau m-N_\sigma\bigl(z^2x^2f_1f_3+x^2y^2f_1f_2\bigr),
\label{eq:N1}\\
N_2&=N_\tau m-N_\sigma\,y^2z^2f_2f_3.
\label{eq:N2}
\end{align}
The factors $16K^3$ and $(mf_1f_2f_3)^{2}$ are symmetric in $x$, $y$,
$z$; a symmetric factor multiplies all three barycentric coordinates
equally and leaves the concurrency locus unchanged (Lemma~6.3,
operation (O5)).  We may therefore replace $f_{131}$ by the center
function
\begin{equation}\label{eq:g131}
g(x,y,z)=\frac{N_1N_2}{S_x},
\end{equation}
which is a rational function of the \emph{squared} side lengths,
homogeneous of degree $30$ in $x$, $y$, $z$.

\emph{Step 2: the concurrence condition as a polynomial identity.}
Apply Theorem~3.1 to the center function \eqref{eq:g131}.  For the
corner triangle $T_A$, let $N_1^{s},N_2^{s},S^{s}$ denote the values of
$N_1,N_2,S_x$ at the ordered side triple $(|FA|,|AE|,|EF|)$, and
$N_1^{t},N_2^{t},S^{t}$ their values at the cyclic order
$(|AE|,|EF|,|FA|)$; define the corresponding quantities for $T_B$ and
$T_C$ analogously.  Then $s_A=N_1^{s}N_2^{s}/S^{s}$,
$t_A=N_1^{t}N_2^{t}/S^{t}$, and cyclically, so $S+T=0$ is equivalent to
\begin{equation}\label{eq:Ncond}
\mathcal N:=
\Bigl(\prod_{V\in\{A,B,C\}}N_1^{s}N_2^{s}\Bigr)
\Bigl(\prod_{V}S^{t}\Bigr)
+
\Bigl(\prod_{V}N_1^{t}N_2^{t}\Bigr)
\Bigl(\prod_{V}S^{s}\Bigr)=0 .
\end{equation}
Because $g$ is homogeneous, multiplying the three squared side lengths
of one corner triangle by a common nonzero factor multiplies $s_V$ and
$t_V$ by the same quantity and does not affect \eqref{eq:Ncond}.  By
Theorem~11.1 we may therefore use the rescaled polynomial triples
\[
\begin{aligned}
\bigl(|FA|^2,\,|AE|^2,\,|EF|^2\bigr)
&\;\propto\;\bigl(c^2(u-w)^2,\;b^2(u-v)^2,\;Q\bigr),\\
\bigl(|DB|^2,\,|BF|^2,\,|FD|^2\bigr)
&\;\propto\;\bigl(a^2(v-u)^2,\;c^2(v-w)^2,\;Q\bigr),\\
\bigl(|EC|^2,\,|CD|^2,\,|DE|^2\bigr)
&\;\propto\;\bigl(b^2(w-v)^2,\;a^2(w-u)^2,\;Q\bigr),
\end{aligned}
\]
where each triple has been divided by
$u^2/\bigl((u-v)(u-w)\bigr)^2$ and its cyclic analogues, and $Q$ is
the quadratic of Theorem~11.1.

\emph{Step 3: substitution of the Euler line.}
The Euler line is
\[
\ell=\bigl[S_A(S_B-S_C):S_B(S_C-S_A):S_C(S_A-S_B)\bigr].
\]
Substituting these values of $u$, $v$, $w$ into the triples above makes
every entry a polynomial in $a^2$, $b^2$, $c^2$ of degree~$5$, and
$\mathcal N$ a homogeneous polynomial in $a^2$, $b^2$, $c^2$ of degree
at most $255$.  A computer algebra computation shows that this
polynomial is identically zero.  (The verification is deterministic:
normalizing $c=1$, the degree of $\mathcal N$ in $a^2$ is at most
$255$; exact evaluation at $260$ distinct values of $a^2$, with $b^2$
kept symbolic, returned the zero polynomial of $\mathbb Q[b^2]$ in
every case, which forces all coefficients of $\mathcal N$ to vanish;
homogeneity then restores general $c$.)

Hence $S+T=0$ whenever the Euler line is an admissible transversal and
the quantities in (3) are defined for $f_{131}$.  By Theorem~3.1, the
cevians $AG$, $BH$, $CI$ are concurrent.
\end{proof}

\end{document}
